\documentclass[letterpaper,12pt]{article}
\setlength{\headheight}{14.49998pt}
\usepackage{fancyhdr}
\usepackage{lipsum,graphicx}
\usepackage{amsmath, amsfonts, amssymb, ragged2e}
\usepackage{amsthm}
\usepackage{bookmark}
\usepackage{bbm}
\usepackage{listings}
\newtheorem{theorem}{Theorem}
\newtheorem{definition}{Definition}

\title{Summary of CPSC 221}
\author{Tom Wang}
\date{Spring, 2024}

\fancypagestyle{plain}{
    \fancyhf{}
    \fancyhead[L]{Tom Wang}
    \fancyhead[R]{\thepage}
}


\begin{document}
\maketitle
\section{Complexity}
A ``good'' algorithm should:\begin{itemize}
    \item Be correct (or at least mostly correct)
    \item finish in a reasonable amount of time (time complexity)
    \item use a reasonable amount of memory (space complexity)
\end{itemize}

\subsection{Time Complexity}
Running time is expressed as a function $T(n):\mathbb{Z}^0 \to \mathbb{R}^0$, where $n$ is the size of the input. We are interested in the asymptotic behavior of $T(n)$ as $n\to \infty$. We use big-O notation to express this.
\subsubsection{Asymptotic Notation}
\begin{definition}

    $T(n)\in O(f(n))$ if there exists $c>0$ and $n_0>0$ such that $T(n)\leq cf(n)$ for all $n\geq n_0$.

    $T(n)\in \Omega(f(n))$ if there exists $c>0$ and $n_0>0$ such that $T(n)\geq cf(n)$ for all $n\geq n_0$.

    $T(n)\in \Theta(f(n))$ if $T(n)\in O(f(n))$ and $T(n)\in \Omega(f(n))$.

    $T(n)\in o(f(n))$ if for all $c>0$, there exists $n_0>0$ such that $T(n)\leq cf(n)$ for all $n\geq n_0$.

    $T(n)\in \omega(f(n))$ if for all $c>0$, there exists $n_0>0$ such that $T(n)\geq cf(n)$ for all $n\geq n_0$.
\end{definition}

The reason we use big-O notation and omit the details is that it is easy to work with. For example, if $T(n)\in O(f(n))$ and $S(n)\in O(g(n))$, then $T(n)+S(n)\in O(f(n)+g(n))$ and $T(n)S(n)\in O(f(n)g(n))$.

We only care about the dominant term in the function (When n gets really large). We also ignore constant factors because we are only interested in the asymptotic behavior.

\subsubsection{Analyze loops}
\begin{lstlisting}[language=C++]
bool hasDuplicate(int arr[], int size){
    for(int i=0; i<size-1; i++){
        for(int j=+1; j<size; j++){
            if(arr[i]==arr[j]){
                return true;
            }
        }
    }
    return false;
}
\end{lstlisting}

Now we analyze the running time of this algorithm. Basically, counts for the worst-case scenario.

It is easy to tell that the outer loop runs $n-1$ times, thus $O(n)$.

The inner loop runs from $n-1$ to $1$ times if we run the loop. However, we can use the fact that the sum of the first $n$ integers is $\frac{n(n-1)}{2}$. So the inner loop runs $\frac{n(n-1)}{2}$ times, thus $O(n^2)$. This is the overall complexity of the nested loop.

Give more examples:\begin{lstlisting}[language=C++]
void candyapple(int n) {
    for (int i = 1; i < n; i *= 3)
    cout << "iteration: " << i << endl;
}
void caramelcorn(int n) {
    for (int i = 0; i * i < 6 * n; i++)
    cout << "iteration: " << i << endl;
}         
\end{lstlisting}
If loop variable is multiplied by a constant, then the loop runs in $O(\log n)$ time. If loop variable is squared, then the loop runs in $O(\sqrt{n})$ time.

So candyapple runs in $O(\log n)$ time and caramelcorn runs in $O(\sqrt{n})$ time.
\begin{proof}

    in candyapple, $i$ is multiplied by 3 each time, so $i=3^k$ after $k$ iterations. So $3^k=n$, so $k=\log_3 n$. So the loop runs $\log_3 n$ times, thus $O(\log n)$.

    in caramelcorn, $i^2=6n$, so $i=\sqrt{6n}$. So the loop runs $\sqrt{6n}$ times, thus $O(\sqrt{n})$.
\end{proof}

lastly, we can analyze the following code:\begin{lstlisting}[language=C++]
int i, j;
for (i = 1; i < 9*n; i = i*2) {
    for (j = n*n; j > 0; j--) {
        ...
    }
}
\end{lstlisting}
The outer loop runs $\log_2 9n$ times, thus $O(\log n)$. The inner loop runs $n^2$ times, thus $O(n^2)$. So the overall complexity is $O(n^2\log n)$.

\end{document}